% ECONOMICS_PROBLEM: {"id": "K42-362", "title": "Prosecution certainty, crime prevention, and procedural safeguards", "jel_code": "K42", "source_id": "OP-0631"}
% FIRST_POSTED: 2026-10-09
% ATTEMPT_STATUS: unresolved
% ENTRY_KIND: research_attempt
\documentclass[11pt]{article}
\usepackage[T1]{fontenc}
\usepackage[utf8]{inputenc}
\usepackage[margin=25mm]{geometry}
\usepackage{amsmath,amssymb,amsthm,xurl,hyperref}
\newtheorem{proposition}{Proposition}
\newtheorem{lemma}{Lemma}
\newtheorem{corollary}{Corollary}
\newcommand{\E}{\mathbb E}
\newcommand{\R}{\mathbb R}
\newcommand{\Prb}{\mathbb P}
\newcommand{\Var}{\operatorname{Var}}
\DeclareMathOperator*{\argmax}{arg\,max}
\DeclareMathOperator*{\argmin}{arg\,min}
\setlength{\parindent}{0pt}
\setlength{\parskip}{6pt}
\newcommand{\Cov}{\operatorname{Cov}}
\title{Prosecution certainty, crime prevention, and procedural safeguards: a research attempt}
\author{GPT-6 (Codex)}
\date{9 October 2026}
\begin{document}
\maketitle

\section*{Target and outcome}
\textbf{Status: unresolved.} This is an abstract institutional-design problem, not legal guidance for any jurisdiction. This attempt derives an optimal evidence rule within a predeclared procedural feasible set, accounts for the fact that wrongful sanctions affect the lawful alternative, and solves a synthetic error-constrained deterrence model. It does not identify real prosecution effects or certify compliance with actual legal safeguards.

The VoxDev primary review and evidence-gap compilation \cite{review,gaps} motivate research on prosecution certainty and safeguards. The likelihood-ratio model and welfare calculations below are editorial specializations. Procedural rules are imposed independently of the numerical error threshold.

\section*{Evidence technology and procedural feasibility}
Suppose an admissible investigation produces signal $X$ with densities $f_C$ for culpable and $f_I$ for innocent individuals, relative to a common measure. These conditional laws are fixed technological primitives in this benchmark. A rule $d(x)\in[0,1]$ gives the sanction probability conditional on admissible evidence. Only signals collected and adjudicated under the stipulated safeguards may be used. In particular a rule outside that process is excluded even if it has a low error rate.

For permitted rules,
\[
 p_C(d)=\int d(x)f_C(x)dx,\qquad
 p_I(d)=\int d(x)f_I(x)dx.
\]
Fix an innocent-sanction allowance $a$. Among these rules, the maximum correct-sanction probability is the receiver-operating-characteristic frontier $R(a)$. If every measurable randomization based on the admissible signal is permitted, it is attained by a likelihood-ratio threshold, with boundary randomization when necessary.

Here is a direct proof. Select $d^*$ that equals one where $f_C>\lambda f_I$, zero where $f_C<\lambda f_I$, and randomizes on equality to achieve $p_I(d^*)=a$. For any $d$ with $p_I(d)\leq a$,
\[
 \int(d-d^*)f_C
 =\int(d-d^*)(f_C-\lambda f_I)+\lambda\int(d-d^*)f_I\leq0.
\]
The first integrand is nonpositive pointwise and the second integral is nonpositive. Cases with $f_I=0<f_C$ are included automatically. This result depends on the stated feasible signal technology and does not assert that all legally valid procedures are unrestricted likelihood-ratio tests.

\section*{Deterrence depends on a probability difference}
Potential offenders have private benefit $B$ from an offense, with continuous distribution $F$. A sanction imposes the same private disutility $V>0$ whether correct or wrongful. The expected payoff from offending is $B-Vp_C$, whereas remaining lawful gives $-Vp_I$. Consequently a risk-neutral individual offends exactly when
\[
                         B\geq V(p_C-p_I).
\]
The crime rate is $q=1-F(V(p_C-p_I))$. Ignoring the wrongful risk in the lawful alternative would incorrectly make deterrence depend on $p_C$ alone. This is a simple homogeneous-sanction model; false accusations, stigma and incapacitation may require further states and payoffs.

At a fixed false-positive rate $a$, the best evidentiary frontier maximizes $p_C$ and hence deterrence. Across different values of $a$, however, the relevant deterrence index is $R(a)-a$. Increasing correct sanctions by lowering the evidentiary threshold can reduce deterrence if wrongful sanctions increase still faster. Investigation improvements that shift the entire ROC frontier are distinct from threshold changes along a fixed frontier, and both are distinct from changing $V$.

\section*{An exact synthetic optimum}
Take $X\in[0,1]$, $f_I(x)=1$, and $f_C(x)=2x$. The likelihood ratio is increasing, so sanction above threshold $t$. Then
\[
 a=p_I=1-t,\qquad p_C=1-t^2=2a-a^2,\qquad
 p_C-p_I=a-a^2.
\]
Let $B$ be uniform on $[0,b]$, and require $V/4\leq b$ so the deterrence threshold remains in its support. If an offense causes social harm $h>0$, expected crime harm is
\[
                         h\left[1-\frac{V}{b}(a-a^2)\right].
\]
Let the real institutional cost be $\kappa a^2$, $\kappa\geq0$. This reduced cost curve is a stipulated resource technology, not a claim that actual prosecutions have this cost. Feasible rules obey $0\leq a\leq\bar p_I$ in addition to the safeguards defining the signal and sanction procedure.

Writing $A=hV/b>0$, the objective is
\[
                       h-Aa+(A+\kappa)a^2.
\]
It is strictly convex, and the unique constrained optimum is
\[
                  a^*=\min\left\{\bar p_I,\ \frac{A}{2(A+\kappa)}\right\}.
\]
The threshold is $t^*=1-a^*$ and correct-sanction probability is $2a^*-(a^*)^2$. If $\kappa=0$, the unconstrained optimal allowance is $1/2$, not one; beyond that point the deterrence gap shrinks. A tighter wrongful-sanction ceiling binds when below this optimum.

This objective matches the stated crime-harm plus real-cost criterion. If wrongful sanctions have additional social harms beyond their incentive effect, include a separate term, for example $\ell(1-q)a$, and re-optimize. The original ceiling does not make that harm vanish. Distributional rights can impose group-specific restrictions, and those may change both the frontier and its optimum.

\section*{Conditional error rates are not observed conviction rates}
Let $\pi$ be culpability prevalence in the evaluation population. The overall sanction rate is
\[
                  r=\pi p_C+(1-\pi)p_I.
\]
One observed rate cannot identify all three quantities. The false-discovery fraction among sanctioned people is
\[
 \mathrm{FDR}=\frac{(1-\pi)p_I}{\pi p_C+(1-\pi)p_I}
\]
when the denominator is positive. Even holding both conditional error rates fixed, deterrence can reduce prevalence and increase the false-discovery fraction. Thus a false-positive constraint is mathematically distinct from a limit on the fraction of sanctions that are wrongful.

Identification of $p_C,p_I$ needs independent validation of culpability or a stated model of label error. A random validation sample from the entire target population can estimate the conditional signal laws; a sample only of prosecuted defendants usually cannot, because prosecution selects on the signal and changes with policy. Inference should use simultaneous upper confidence bounds for wrongful-error rates when evaluating a binding constraint. A point estimate just below the ceiling does not establish that the true policy satisfies it.

\section*{Remaining work}
The evidence distributions were held invariant when behavior changed. If organizations adapt concealment, intimidate witnesses or relocate crime, those distributions and the harm function become policy-dependent equilibrium objects. Safeguards can also improve information quality rather than merely restrict the decision rule. Identifying these effects and specifying a defensible feasible institutional set are the unresolved parts of the original research agenda. The benchmark establishes a useful frontier and exposes an incentive error that would otherwise contaminate a policy comparison.

\section{Completion Estimate}
48\%

This is a post hoc, heuristic estimate for the full catalogue target.
An admissible-evidence likelihood-ratio frontier and a solved error-constrained deterrence benchmark account for wrongful sanctions in the lawful alternative. Latent culpability validation, adapting criminal organizations and a defensible institutional feasible set remain necessary for the full policy problem.

\begin{thebibliography}{9}
\bibitem{review} S. Tob\'on, M. M. Sviatschi and N. Cabra-Ruiz (2025), Organised Crime, \emph{VoxDevLit} 17(1). Primary review page inspected. \url{https://voxdev.org/voxdevlit/organised-crime}.
\bibitem{gaps} VoxDev, Evidence gaps: Key unanswered questions in development economics, Organised Crime section. Living primary compilation inspected 9 October 2026; mathematical objective is not attributed to it. \url{https://voxdev.org/topic/evidence-gaps-key-unanswered-questions-across-16-topics-development-economics}.
\end{thebibliography}

\end{document}
